\documentclass[12pt,a4paper]{article}
\usepackage[utf8]{inputenc}
\usepackage[T1]{fontenc}
\usepackage[english,romanian]{babel}
\usepackage{graphicx}
\usepackage{subcaption}
\usepackage{amsmath,amssymb}
\usepackage{fancyhdr}
\usepackage[a4paper,margin=2.5cm]{geometry}
\usepackage{dirtytalk}
\usepackage{csquotes}
\usepackage[hidelinks]{hyperref}
\usepackage{listings}
\usepackage{xcolor}
\usepackage{array}
\usepackage{booktabs}

\lstset{
    language=Python,
    basicstyle=\ttfamily\small,
    numbers=left,
    numberstyle=\tiny,
    frame=single,
    breaklines=true,
    showstringspaces=false,
    captionpos=b
}
% listings nu înțelege direct diacriticele UTF-8, de aceea le traducem explicit.
\lstset{
    literate=
        {Ă}{{\u{A}}}1 {ă}{{\u{a}}}1
        {Â}{{\^A}}1 {â}{{\^a}}1
        {Î}{{\^I}}1 {î}{{\^i}}1
        {Ș}{{\textcommabelow{S}}}1 {ș}{{\textcommabelow{s}}}1
        {Ț}{{\textcommabelow{T}}}1 {ț}{{\textcommabelow{t}}}1
}
\lstdefinestyle{terminal}{
    language={},
    numbers=none,
    backgroundcolor=\color{black!5},
    basicstyle=\ttfamily\footnotesize
}
\setlength{\parindent}{1.25cm}
\setlength{\parskip}{0.25em}

\begin{document}
\begin{titlepage}
\begin{center}
    % Înlocuiți calea dacă logoul este stocat în altă parte.
    \includegraphics[width=0.3\textwidth]{lab1/image .png}\\[0.2cm]
    \textbf{MINISTERUL EDUCAȚIEI, CULTURII ȘI CERCETĂRII\\
    AL REPUBLICII MOLDOVA}\\[0.3cm]
    \textbf{Universitatea Tehnică a Moldovei\\
    Facultatea Calculatoare, Informatică și Microelectronică\\
    Departamentul Inginerie Software și Automatică}\\[2cm]
    \Huge\textbf{Raport}\\[0.5cm]
    \Large Lucrarea de laborator nr. 4\\[0.3cm]
    \large\textbf{la Criptografie și Securitate}\\[0.8cm]
    \textbf{Ciur Alexandrina, FAF-241}\\[2.5cm]
    \begin{flushright}
        \textit{Verificat de:}\\
        \textbf{Zaica M.}\\
        DISA, FCIM, UTM
    \end{flushright}
    \vfill
    Chișinău -- 2026
\end{center}
\end{titlepage}

\section*{Tema: Cifruri bloc. Algoritmul DES}

\section{Scopul lucrării}

Studierea cifrurilor bloc și a algoritmului DES (Data Encryption Standard) și
implementarea unui element al acestuia -- aplicarea cutiilor de substituție
(cutiile S) asupra unui bloc de 48 de biți, cu afișarea tabelelor folosite și a
tuturor pașilor intermediari.

\section{Sarcina lucrării}

Numărul sarcinii se alege după formula $nr\_sarcina = n \bmod 11$, unde $n$ este
numărul de ordine al studentului în lista grupei. Pentru $n = 7$ obținem
$7 \bmod 11 = 7$, prin urmare a fost realizată \textbf{Sarcina 2.7}:

\begin{quote}
\itshape
În algoritmul DES este dat $B_1B_2B_3B_4B_5B_6B_7B_8$. De aflat
\[
    S_1(B_1)S_2(B_2)S_3(B_3)S_4(B_4)S_5(B_5)S_6(B_6)S_7(B_7)S_8(B_8).
\]
Să fie afișate la ecran tabelele utilizate și toți pașii intermediari. Datele de
intrare să fie posibil de introdus de utilizator sau de generat în mod aleatoriu.
\end{quote}

\section{Considerații teoretice}

\subsection{Cifrurile bloc}

Un \textbf{cifru bloc} împarte mesajul în blocuri de lungime fixă și criptează
fiecare bloc separat, folosind aceeași cheie. Pentru un bloc de $n$ biți, cifrul
este o funcție care, pentru o cheie fixată, transformă fiecare dintre cele $2^n$
blocuri posibile ale textului clar într-un bloc unic al criptogramei. Dacă
ultimul bloc al mesajului este mai scurt, el se completează (\textit{padding}).

Majoritatea cifrurilor bloc moderne se bazează pe două principii propuse de
Claude Shannon:
\begin{itemize}
    \item \textbf{confuzia} -- legătura dintre cheie și criptogramă trebuie să fie
    cât mai complicată (se obține prin \textbf{substituții});
    \item \textbf{difuzia} -- fiecare bit al textului clar trebuie să influențeze
    cât mai mulți biți ai criptogramei (se obține prin \textbf{permutări}).
\end{itemize}
Aceste operații se aplică de mai multe ori, în \textbf{runde}, fiecare rundă
folosind o cheie de rundă diferită, obținută din cheia principală.

\subsection{Scurt istoric al algoritmului DES}

DES a fost elaborat de compania IBM (pe baza cifrului Lucifer) și a fost adoptat
în 1977 ca standard federal al SUA (FIPS 46). Timp de peste 20 de ani a fost cel
mai folosit cifru bloc din lume. Cheia efectivă de doar 56 de biți a devenit însă
prea scurtă: în 1998 mașina \enquote{Deep Crack} a găsit o cheie DES prin căutare
exhaustivă în mai puțin de 3 zile. De aceea, standardul FIPS 46-3 (1999)
recomandă varianta \textbf{Triple DES} (3DES), iar în 2001 DES a fost înlocuit de
algoritmul AES.

\subsection{Structura generală a algoritmului DES}

DES criptează blocuri de \textbf{64 de biți} cu o cheie de \textbf{64 de biți}, din
care doar 56 sunt folosiți efectiv (fiecare al 8-lea bit este bit de paritate).
Algoritmul are o \textbf{structură Feistel} cu 16 runde:

\begin{enumerate}
    \item asupra blocului de 64 de biți se aplică \textbf{permutarea inițială}
    $IP$, iar rezultatul se împarte în două jumătăți de câte 32 de biți:
    $L_0$ (stânga) și $R_0$ (dreapta);
    \item se execută 16 runde; în runda $i$ ($i = 1, \dots, 16$):
    \[
        L_i = R_{i-1}, \qquad R_i = L_{i-1} \oplus f(R_{i-1}, K_i),
    \]
    unde $\oplus$ este operația XOR (adunare modulo 2), iar $K_i$ este cheia de
    rundă de 48 de biți;
    \item după runda 16, jumătățile se unesc în ordine inversă, $R_{16}L_{16}$, și
    se aplică \textbf{permutarea finală} $IP^{-1}$, obținându-se criptograma de
    64 de biți.
\end{enumerate}

Decriptarea folosește același algoritm, dar cheile de rundă se aplică în ordine
inversă ($K_{16}, K_{15}, \dots, K_1$). Aceasta este o proprietate a structurii
Feistel: funcția $f$ nu trebuie să fie inversabilă.

\subsection{Generarea cheilor de rundă}

\begin{enumerate}
    \item Din cheia de 64 de biți se aplică permutarea \textbf{PC-1}, care elimină
    biții de paritate și dă cheia $K^+$ de 56 de biți.
    \item $K^+$ se împarte în două jumătăți de câte 28 de biți: $C_0$ și $D_0$.
    \item Pentru $i = 1, \dots, 16$, $C_i$ și $D_i$ se obțin din $C_{i-1}$ și
    $D_{i-1}$ prin \textbf{deplasare ciclică la stânga} cu 1 sau 2 poziții, conform
    tabelului:
    \begin{center}
    \small
    \begin{tabular}{|c|*{16}{c|}}
    \hline
    Runda $i$ & 1 & 2 & 3 & 4 & 5 & 6 & 7 & 8 & 9 & 10 & 11 & 12 & 13 & 14 & 15 & 16 \\ \hline
    Deplasări & 1 & 1 & 2 & 2 & 2 & 2 & 2 & 2 & 1 & 2 & 2 & 2 & 2 & 2 & 2 & 1 \\ \hline
    \end{tabular}
    \end{center}
    \item Cheia de rundă $K_i$ de 48 de biți se obține aplicând permutarea
    \textbf{PC-2} asupra concatenării $C_iD_i$.
\end{enumerate}

\subsection{Funcția $f$}

Funcția $f(R_{i-1}, K_i)$ primește 32 de biți și întoarce 32 de biți:
\begin{enumerate}
    \item \textbf{expansiunea} $E$ transformă $R_{i-1}$ din 32 în 48 de biți (unii
    biți se repetă);
    \item rezultatul se adună modulo 2 cu cheia de rundă:
    $K_i \oplus E(R_{i-1}) = B_1B_2B_3B_4B_5B_6B_7B_8$, unde fiecare $B_j$ are
    6 biți;
    \item fiecare $B_j$ trece prin cutia $S_j$ și se obțin 4 biți, deci în total
    $8 \cdot 4 = 32$ de biți: $S_1(B_1)S_2(B_2)\dots S_8(B_8)$ -- \textbf{pasul
    realizat în această lucrare};
    \item asupra celor 32 de biți se aplică permutarea $P$:
    \[
        f(R_{i-1}, K_i) = P\bigl(S_1(B_1)S_2(B_2)\dots S_8(B_8)\bigr).
    \]
\end{enumerate}

\subsection{Cutiile S}

DES are 8 cutii S ($S_1, \dots, S_8$). Fiecare este un tabel cu \textbf{4 rânduri}
(numerotate 0--3) și \textbf{16 coloane} (numerotate 0--15), cu valori între 0 și
15. Pentru un bloc de 6 biți $B_j = b_1b_2b_3b_4b_5b_6$:
\begin{itemize}
    \item \textbf{rândul} este numărul scris în binar de primul și ultimul bit:
    $r = (b_1b_6)_2$;
    \item \textbf{coloana} este numărul scris în binar de cei 4 biți din mijloc:
    $c = (b_2b_3b_4b_5)_2$;
    \item $S_j(B_j)$ este valoarea $S_j[r][c]$, scrisă pe \textbf{4 biți}.
\end{itemize}

Cutiile S sunt singura componentă \textbf{neliniară} a algoritmului DES: toate
celelalte operații (permutări, expansiune, XOR) sunt liniare. Fără ele, DES ar
putea fi spart prin rezolvarea unui sistem de ecuații liniare. Fiecare rând al
fiecărei cutii conține o singură dată toate numerele 0--15, iar modificarea unui
singur bit de intrare schimbă cel puțin doi biți de ieșire, ceea ce asigură
difuzia de la o rundă la alta.

Tabelul cutiei $S_1$ (celelalte șapte au aceeași formă și sunt afișate de program):
\begin{center}
\small
\renewcommand{\arraystretch}{1.2}
\setlength{\tabcolsep}{4pt}
\begin{tabular}{|c|*{16}{c|}}
\hline
 & 0 & 1 & 2 & 3 & 4 & 5 & 6 & 7 & 8 & 9 & 10 & 11 & 12 & 13 & 14 & 15 \\ \hline
0 & 14 & 4 & 13 & 1 & 2 & 15 & 11 & 8 & 3 & 10 & 6 & 12 & \textbf{5} & 9 & 0 & 7 \\ \hline
1 & 0 & 15 & 7 & 4 & 14 & 2 & 13 & 1 & 10 & 6 & 12 & 11 & 9 & 5 & 3 & 8 \\ \hline
2 & 4 & 1 & 14 & 8 & 13 & 6 & 2 & 11 & 15 & 12 & 9 & 7 & 3 & 10 & 5 & 0 \\ \hline
3 & 15 & 12 & 8 & 2 & 4 & 9 & 1 & 7 & 5 & 11 & 3 & 14 & 10 & 0 & 6 & 13 \\ \hline
\end{tabular}
\end{center}

\subsection{Exemplu de calcul manual}

Fie $B_1 = \texttt{011000}$:
\begin{itemize}
    \item primul și ultimul bit: $\texttt{0}\dots\texttt{0} \rightarrow
    (00)_2 = 0$, deci rândul $r = 0$;
    \item biții din mijloc: $(1100)_2 = 12$, deci coloana $c = 12$;
    \item $S_1[0][12] = 5$ (evidențiat în tabelul de mai sus), deci
    $S_1(B_1) = 5 = \texttt{0101}$.
\end{itemize}

Aplicând aceeași regulă tuturor celor 8 blocuri din exemplul descris în
materialul \enquote{The DES Algorithm Illustrated}, obținem:

\begin{center}
\renewcommand{\arraystretch}{1.2}
\begin{tabular}{@{}cccccc@{}}
\toprule
$j$ & $B_j$ & \textbf{Rând} & \textbf{Coloană} & $S_j[r][c]$ & $S_j(B_j)$ \\
\midrule
1 & \texttt{011000} & 0 & 12 & 5  & \texttt{0101} \\
2 & \texttt{010001} & 1 & 8  & 12 & \texttt{1100} \\
3 & \texttt{011110} & 0 & 15 & 8  & \texttt{1000} \\
4 & \texttt{111010} & 2 & 13 & 2  & \texttt{0010} \\
5 & \texttt{100001} & 3 & 0  & 11 & \texttt{1011} \\
6 & \texttt{100110} & 2 & 3  & 5  & \texttt{0101} \\
7 & \texttt{010100} & 0 & 10 & 9  & \texttt{1001} \\
8 & \texttt{100111} & 3 & 3  & 7  & \texttt{0111} \\
\bottomrule
\end{tabular}
\end{center}

Rezultatul final (32 de biți) este
$\texttt{0101 1100 1000 0010 1011 0101 1001 0111}$, adică
$\texttt{5C82B597}$ în hexazecimal.

\section{Implementarea}

Programul a fost realizat în limbajul \textbf{Python 3} (fișierul
\texttt{lab4/lab4.py}). Mai jos sunt prezentate părțile principale ale codului.

\subsection{Constantele și tabelele cutiilor S}

\begin{lstlisting}[caption={Parametrii și începutul tabelelor S (FIPS 46-3)}]
INPUT_BITS = 48
BLOCK_COUNT = 8
BLOCK_BITS = 6
OUTPUT_BITS = 4

# Cutiile S1-S8 din standardul FIPS 46-3: 4 randuri x 16 coloane fiecare.
S_BOXES = (
    (
        (14, 4, 13, 1, 2, 15, 11, 8, 3, 10, 6, 12, 5, 9, 0, 7),
        (0, 15, 7, 4, 14, 2, 13, 1, 10, 6, 12, 11, 9, 5, 3, 8),
        (4, 1, 14, 8, 13, 6, 2, 11, 15, 12, 9, 7, 3, 10, 5, 0),
        (15, 12, 8, 2, 4, 9, 1, 7, 5, 11, 3, 14, 10, 0, 6, 13),
    ),
    # ... S2 - S8 au aceeasi forma
)
\end{lstlisting}

\texttt{S\_BOXES} conține cele 8 cutii; \texttt{S\_BOXES[0]} este $S_1$,
\texttt{S\_BOXES[7]} este $S_8$. Fiecare cutie este formată din 4 rânduri a câte
16 valori, copiate din standardul FIPS 46-3.

\subsection{Validarea și generarea datelor de intrare}

\begin{lstlisting}[caption={Verificarea celor 48 de biți și generarea aleatorie}]
def validate_bits(raw):
    """Remove spaces and check the input, returning (bits, error_message)."""
    bits = raw.replace(" ", "")
    for character in bits:
        if character not in "01":
            return None, f"Caracter invalid: {character!r}. Folositi doar 0 si 1."
    if len(bits) != INPUT_BITS:
        return None, (f"Trebuie exact {INPUT_BITS} de biti "
                      f"(ati introdus {len(bits)}).")
    return bits, None


def random_bits():
    """Generate a random 48-bit string."""
    return "".join(random.choice("01") for _ in range(INPUT_BITS))
\end{lstlisting}

Utilizatorul poate scrie biții cu spații între grupe (de exemplu, câte 6),
deoarece spațiile sunt eliminate. Programul verifică faptul că sunt folosite doar
cifrele 0 și 1 și că sunt exact 48 de biți. A doua funcție generează aleatoriu un
șir de 48 de biți.

\subsection{Aplicarea cutiilor S}

\begin{lstlisting}[caption={Împărțirea în blocuri și căutarea în cutiile S}]
def split_blocks(bits):
    """Split 48 bits into the eight 6-bit blocks B1..B8."""
    return [bits[index:index + BLOCK_BITS] for index in range(0, len(bits), BLOCK_BITS)]


def sbox_lookup(box_number, block):
    """Apply S-box number box_number (1-8) to a 6-bit block.

    Returns (row, column, value, output_bits).
    """
    row = int(block[0] + block[-1], 2)
    column = int(block[1:-1], 2)
    value = S_BOXES[box_number - 1][row][column]
    return row, column, value, format(value, f"0{OUTPUT_BITS}b")


def apply_sboxes(bits):
    """Return S1(B1)S2(B2)...S8(B8) as a 32-bit string."""
    return "".join(
        sbox_lookup(number, block)[3]
        for number, block in enumerate(split_blocks(bits), start=1)
    )
\end{lstlisting}

Funcția \texttt{sbox\_lookup} realizează exact regula descrisă în partea teoretică:
\texttt{block[0] + block[-1]} formează rândul din primul și ultimul bit,
\texttt{block[1:-1]} sunt cei 4 biți din mijloc (coloana), iar \texttt{int(..., 2)}
transformă numărul din binar în zecimal. Valoarea găsită este scrisă înapoi pe
4 biți cu \texttt{format(value, "04b")}. Funcția \texttt{apply\_sboxes} aplică
fiecărui bloc $B_j$ cutia $S_j$ și lipește cele 8 rezultate.

\subsection{Afișarea pașilor intermediari}

\begin{lstlisting}[caption={Afișarea fiecărui pas și a rezultatului final}]
def _run_sboxes(bits):
    print("\nIntrarea (48 de biti):", _in_groups(bits, BLOCK_BITS))
    for number, block in enumerate(split_blocks(bits), start=1):
        row, column, value, output = sbox_lookup(number, block)
        print(f"\n--- B{number} = {block} ---")
        print(f"Randul:  primul si ultimul bit {block[0]}{block[-1]} -> {row}")
        print(f"Coloana: bitii din mijloc {block[1:-1]} -> {column}")
        _show_sbox(number, row, column)
        print(f"S{number}(B{number}) = S{number}[{row}][{column}] = {value} -> {output}")

    result = apply_sboxes(bits)
    print("\nS1(B1)S2(B2)...S8(B8) (32 de biti):", _in_groups(result, OUTPUT_BITS))
    print("In hexazecimal:", format(int(result, 2), "08X"))
\end{lstlisting}

Pentru fiecare bloc se afișează rândul, coloana, tabelul cutiei folosite (cu
valoarea aleasă pusă între paranteze drepte, prin funcția \texttt{\_show\_sbox}) și
rezultatul de 4 biți. La final se afișează cei 32 de biți obținuți, în binar și în
hexazecimal.

\section{Rezultatele execuției}

Programul se lansează din folderul \texttt{lab4} cu comanda:
\begin{lstlisting}[style=terminal]
(.venv) PS D:\Criptografia\lab4> py lab4.py
\end{lstlisting}

\subsection{Introducerea datelor de către utilizator}

Utilizatorul alege opțiunea \texttt{1} și introduce cei 48 de biți. Pentru
fiecare bloc programul afișează pașii și tabelul folosit. Mai jos este prezentat
complet primul bloc; pentru blocurile $B_2$--$B_8$ s-au păstrat doar rezultatele,
iar tabelele au fost omise din raport.

\begin{lstlisting}[style=terminal, basicstyle=\ttfamily\scriptsize, caption={Aplicarea cutiilor S pentru datele introduse de utilizator}]
=== DES: cutiile S (sarcina 2.7) ===
1. Introduc B1B2...B8
2. Generez B1B2...B8 aleatoriu
0. Iesire
Alege (0-2): 1
...
B1B2...B8 (48 de biti): 000011110101010100001111101010011001101001101010

Intrarea (48 de biti): 000011 110101 010100 001111 101010 011001 101001 101010

--- B1 = 000011 ---
Randul:  primul si ultimul bit 01 -> 1
Coloana: bitii din mijloc 0001 -> 1
Tabelul S1 (valoarea aleasa este intre [ ]):
          0    1    2    3    4    5    6    7    8    9   10   11   12   13   14   15
  r0:    14    4   13    1    2   15   11    8    3   10    6   12    5    9    0    7
  r1:     0 [15]    7    4   14    2   13    1   10    6   12   11    9    5    3    8
  r2:     4    1   14    8   13    6    2   11   15   12    9    7    3   10    5    0
  r3:    15   12    8    2    4    9    1    7    5   11    3   14   10    0    6   13
S1(B1) = S1[1][1] = 15 -> 1111
...
S2(B2) = S2[3][10] = 7 -> 0111
...
S3(B3) = S3[0][10] = 12 -> 1100
...
S4(B4) = S4[1][7] = 3 -> 0011
...
S5(B5) = S5[2][5] = 13 -> 1101
...
S6(B6) = S6[1][12] = 0 -> 0000
...
S7(B7) = S7[3][4] = 1 -> 0001
...
S8(B8) = S8[2][5] = 12 -> 1100

S1(B1)S2(B2)...S8(B8) (32 de biti): 1111 0111 1100 0011 1101 0000 0001 1100
In hexazecimal: F7C3D01C
\end{lstlisting}

Rezultatul poate fi verificat manual. De exemplu, pentru $B_1 = \texttt{000011}$
primul și ultimul bit dau rândul $(01)_2 = 1$, biții din mijloc dau coloana
$(0001)_2 = 1$, iar în tabelul $S_1$ găsim $S_1[1][1] = 15 = \texttt{1111}$.
De asemenea, pentru exemplul din partea teoretică programul afișează rezultatul
\texttt{5C82B597}, care coincide cu cel din materialul
\enquote{The DES Algorithm Illustrated}. Aceasta confirmă corectitudinea tabelelor
și a implementării.

% Pentru a adăuga o captură de ecran, încărcați imaginea în Overleaf și
% decomentați blocul de mai jos (schimbați numele fișierului dacă e altul).
% \begin{figure}[h!]
%     \centering
%     \includegraphics[width=0.9\textwidth]{lab4/introducere.png}
%     \caption{Aplicarea cutiilor S în terminal}
% \end{figure}

\subsection{Generarea aleatorie a datelor}

La alegerea opțiunii \texttt{2}, programul generează singur cei 48 de biți și
afișează aceiași pași.

\begin{lstlisting}[style=terminal, caption={Aplicarea cutiilor S pentru date generate aleatoriu}]
Alege (0-2): 2

Intrarea (48 de biti): 110011 001000 001100 110000 101111 101011 111110 001101

--- B1 = 110011 ---
Randul:  primul si ultimul bit 11 -> 3
Coloana: bitii din mijloc 1001 -> 9
...
S1(B1) = S1[3][9] = 11 -> 1011
...
S2(B2) = S2[0][4] = 6 -> 0110
...
S3(B3) = S3[0][6] = 15 -> 1111
...
S4(B4) = S4[2][8] = 15 -> 1111
...
S5(B5) = S5[3][7] = 13 -> 1101
...
S6(B6) = S6[3][5] = 5 -> 0101
...
S7(B7) = S7[2][15] = 2 -> 0010
...
S8(B8) = S8[1][6] = 7 -> 0111

S1(B1)S2(B2)...S8(B8) (32 de biti): 1011 0110 1111 1111 1101 0101 0010 0111
In hexazecimal: B6FFD527
\end{lstlisting}

% \begin{figure}[h!]
%     \centering
%     \includegraphics[width=0.9\textwidth]{lab4/aleatoriu.png}
%     \caption{Generarea aleatorie a datelor de intrare}
% \end{figure}

\subsection{Tratarea datelor incorecte}

Dacă utilizatorul introduce alte caractere decât 0 și 1 sau un număr greșit de
biți, programul afișează un mesaj explicativ și cere datele din nou. În același
mod este tratată și alegerea unei opțiuni inexistente în meniu.

\begin{lstlisting}[style=terminal, caption={Mesaje de eroare la introducerea datelor greșite}]
Alege (0-2): 1
B1B2...B8 (48 de biti): b1b3b45bbkdjh;
Caracter invalid: 'b'. Folositi doar 0 si 1.
B1B2...B8 (48 de biti): 0101010000000000011111111111111
Trebuie exact 48 de biti (ati introdus 31).
B1B2...B8 (48 de biti): 000011110101010100001111101010011001101001101010
\end{lstlisting}

\section{Concluzii}

În cadrul acestei lucrări de laborator au fost studiate cifrurile bloc și
algoritmul DES: structura Feistel cu 16 runde, generarea celor 16 chei de rundă
din cheia principală și funcția $f$. A fost implementat în limbajul Python pasul
de aplicare a cutiilor S, care transformă cei 48 de biți
$B_1B_2\dots B_8 = K_i \oplus E(R_{i-1})$ în cei 32 de biți
$S_1(B_1)S_2(B_2)\dots S_8(B_8)$. Programul permite introducerea datelor de către
utilizator sau generarea lor aleatorie, verifică datele introduse și afișează
pentru fiecare bloc rândul, coloana, tabelul folosit și rezultatul.

Am observat că fiecare cutie S reduce 6 biți la 4 biți: primul și ultimul bit
aleg unul dintre cele 4 rânduri, iar cei 4 biți din mijloc aleg coloana. Cutiile S
sunt singura parte neliniară a algoritmului DES și de ele depinde în cea mai mare
măsură securitatea acestuia. Corectitudinea implementării a fost verificată pe
exemplul din materialul \enquote{The DES Algorithm Illustrated}, pentru care s-a
obținut rezultatul așteptat \texttt{5C82B597}.

Deși structura algoritmului DES este bine gândită, cheia efectivă de 56 de biți
este astăzi prea scurtă și poate fi găsită prin căutare exhaustivă. De aceea DES a
fost înlocuit de Triple DES și apoi de AES, însă ideile sale (rundele, structura
Feistel, substituțiile și permutările) se regăsesc și în cifrurile bloc moderne.

Codul sursă al lucrării este disponibil în repozitoriul GitHub:\\
\url{https://github.com/alecsyyyk/Cryptography-and-Security}

\end{document}
